% Evaluación de <ASIGNATURA> — <GRADO> — <PERIODO> <AÑO>
% Compilar: bash .claude/skills/evaluaciones-latex/scripts/compilar.sh <este archivo>
% (genera la versión del estudiante y la del docente con las respuestas marcadas)
\documentclass[carta,doscolumnas,compacta]{../plantilla/evaluacion}
\usepackage{../plantilla/evaluacion-extras}  % fragmentos, retratos, líneas de tiempo, historietas...

% Colegio, logo, color y docente: de curriculo/config/colegio.yaml y docente.yaml.
% Para otro valor en este documento: \colegio{…} \logo{archivo.png} \colorcolegio{166534} \docente{…}

\asignatura{<Asignatura>}
\grado{<Sexto, Séptimo, ...>}
\periodo{<Cuarto periodo>}
\anio{<2026>}
\fecha{<dd/mm/aaaa>}
\titulo{Evaluación de <Asignatura>}
\instrucciones{Tiempo: 1 hora. Todo lo que necesitas para responder está en los textos, tablas,
esquemas e imágenes: léelos con atención antes de contestar. En las preguntas de selección marca
una sola respuesta (A, B, C o D). En las de \textbf{afirmación y razón} decide si cada parte es
verdadera y si la razón explica la afirmación. En las abiertas se califica la claridad y la solidez
de tus argumentos.}

% \afirmacionrazon, \afirmacion, afirmaciones, \fragmento... vienen de evaluacion-extras.

\begin{document}
\encabezado

\begin{preguntas}

% ---------------------------------------------------------------- Bloque 1
\begin{contexto}[<Título del texto o caso>]
\grupo{Responde las preguntas 1 a 3 con base en el texto.}
<Resumen de la película, lectura, caso o teoría: todo lo necesario para responder.>
\fragmento{<Texto citado>}{<Autor, \emph{Obra} libro, pasaje>, adaptado}
% \retrato[0.30]{img/<carpeta>/<archivo>.jpg}{<Pie: autor, obra, fecha>}{<texto al lado>}
\end{contexto}

% Selección única (la letra entre corchetes es la sorteada)
\begin{pregunta}[B]
<Enunciado que pide analizar, inferir o aplicar.>
\opcioneslargas{<A>}{<B: correcta>}{<C>}{<D>}
\justificacion{<Por qué B es correcta y qué error revela cada distractor.>}
\especificacion{<Interpretativa | Argumentativa | Propositiva>}{<Evidencia que se evalúa>}
\end{pregunta}

% Afirmaciones I, II, III
\begin{pregunta}[C]
Lee las afirmaciones:
\begin{afirmaciones}
\item <...>
\item <...>
\item <...>
\end{afirmaciones}
¿Cuáles son correctas?
\opciones{Solo I}{Solo III}{I y II}{I, II y III}
\end{pregunta}

% Afirmación y razón
\begin{pregunta}[A]
\afirmacion{<afirmación>}{<razón>}
\afirmacionrazon
\end{pregunta}

% Opciones que son imágenes (TikZ o \includegraphics)
\begin{pregunta}[D]
<Enunciado>
\opcionesdiagrama{<imagen A>}{<imagen B>}{<imagen C>}{<imagen D: correcta>}
\end{pregunta}

% ---------------------------------------------------------------- Abiertas
\begin{pregunta}[abierta]
<Pregunta que pida argumentar, decidir o diseñar.>
\abierta[8]
\modelo{<Respuesta esperada y criterios de calificación.>}
\end{pregunta}

\end{preguntas}

\autoevaluacion
\creditos{<Fuentes de los fragmentos e imágenes con autor y licencia.>}
\tablaespecificaciones
\imprimirclave
\end{document}
